\chapter{यूक्लिडीय समष्टियाँ}\label{ch:b1:euclid}

किसी वास्तविक \omterm{def:b1:vspaces:def}{सदिश समष्टि} में \omterm{def:b1:euclid:def}{अंतर्गुणन} जोड़ देने से ज्यामितीय अवधारणाएँ
ख़रीद ली जाती हैं --- लंबाइयाँ, कोण, लांबिकता, दूरियाँ --- और साथ ही एक ऐसी
प्रमेय जो पूरे अध्याय पर मीनार की तरह खड़ी है: हर \omterm{def:b1:vspaces:subspace}{उपसमष्टि} का एक \omterm{thm:b1:euclid:projection}{लांबिक
प्रक्षेप} होता है, जो ग्राम--श्मिट से संगणनीय है और सबसे छोटी दूरी साकार करता
है। समतल के \omterm{def:b1:euclid:isometry}{समदूरक} अध्याय और वर्ष की ज्यामिति, दोनों को बंद करते हैं।

सर्वत्र $E$ एक \emph{वास्तविक} \omterm{def:b1:vspaces:def}{सदिश समष्टि} है।

\section{अंतर्गुणन}

\begin{definition}\label{def:b1:euclid:def}
$E$ पर \emph{अंतर्गुणन}\index{अंतर्गुणन} ऐसा \omterm{def:b1:logic:map}{प्रतिचित्रण}
$\langle\cdot,\cdot\rangle \colon E \times E \to \R$ है जो द्विरैखिक, सममित
और धन-निश्चित हो ($x \neq 0$ के लिए $\langle x, x\rangle > 0$)। इस प्रकार
सज्जित \omterm{def:b1:findim:def}{परिमित-विमीय} समष्टि \emph{यूक्लिडीय समष्टि}\index{यूक्लिडीय समष्टि}
कहलाती है। $x$ का \emph{मानक} $\norm{x} = \sqrt{\langle x, x\rangle}$ है, और
$d(x, y) = \norm{x - y}$।
\end{definition}

\begin{example}\label{ex:b1:euclid:examples}
$\R^n$ पर: विहित गुणन $\langle x, y\rangle = \sum x_i y_i$।
$C(\intcc{a}{b})$ पर: $\langle f, g \rangle = \int_a^b fg$ (धन-निश्चितता
\cref{thm:b1:integration:props}~(4) है)। $\R_n[X]$ पर: $\langle P, Q\rangle
= \int_0^1 PQ$, अथवा $n+1$ भिन्न बिंदुओं पर $\sum_{i} P(x_i)Q(x_i)$।
\end{example}

\begin{example}[दो बहुपदों के बीच का कोण]
\label{ex:b1:euclid:polyangle}
\omterm{def:b1:euclid:def}{अंतर्गुणन} चुन लेने के बाद \emph{किन्हीं} भी दो अशून्य सदिशों के बीच कोण होता
है, $\cos\theta = \frac{\langle x, y\rangle}{\norm x\,\norm y}$ के द्वारा
(कोशी--श्वार्ट्ज़ से यह वैध कोज्या है)। $\int_0^1$ में $X$ और $X^2$ के लिए:
\[
\langle X, X^2\rangle = \frac14, \qquad
\norm X = \frac1{\sqrt3}, \qquad \norm{X^2} = \frac1{\sqrt5},
\qquad
\cos\theta = \frac{1/4}{1/\sqrt{15}} = \frac{\sqrt{15}}{4}
\approx 0.968 :
\]
अर्थात् लगभग $14.5$ अंश का कोण --- $\intcc{0}{1}$ पर $x$ और $x^2$ के आलेख
वर्ग-माध्य के अर्थ में ``लगभग समांतर'' हैं, और इसीलिए उस साझा दिशा को हटा
देने पर (नीचे ग्राम--श्मिट) केवल छोटा-सा सुधार $X^2 - X + \frac16$ बचता है।
\end{example}

\begin{theorem}[कोशी--श्वार्ट्ज़; मानक के गुण]\label{thm:b1:euclid:cs}
सभी $x, y \in E$ के लिए:\index{कोशी--श्वार्ट्ज़ असमिका}
\[
\abs{\langle x, y\rangle} \leq \norm x\, \norm y ,
\]
और समानता तभी होती है जब $x, y$ अनुपाती हों। परिणामस्वरूप $\norm\cdot$
त्रिभुज असमिका $\norm{x + y} \leq \norm x + \norm y$ संतुष्ट करता है (तथा
$\norm{\lambda x} = \abs\lambda \norm x$, $\norm x = 0 \iff x = 0$)। इसके
अतिरिक्त:
\[
\norm{x+y}^2 = \norm x^2 + 2\langle x, y\rangle + \norm y^2,
\qquad
\langle x, y \rangle = \tfrac14\bigl(\norm{x+y}^2 -
\norm{x-y}^2\bigr) .
\]
\end{theorem}

\begin{proof}
यदि $y = 0$, तो सब कुछ तुच्छ है। अन्यथा द्विघातीय $t \mapsto \norm{x + ty}^2
= \norm x^2 + 2t\langle x, y\rangle + t^2 \norm y^2$ सभी $t$ के लिए $\geq 0$
है: उसका विविक्तकर $\leq 0$ है, और यही कोशी--श्वार्ट्ज़ है; समानता का अर्थ
है कोई दुहरा मूल $t_0$, अर्थात् $x + t_0 y = 0$ (निश्चितता): यानी अनुपातिता।
त्रिभुज असमिका: प्रसार कीजिए,
\[
\norm{x + y}^2 = \norm x^2 + 2\langle x, y\rangle + \norm y^2
\leq \norm x^2 + 2\norm x\,\norm y + \norm y^2
= \bigl(\norm x + \norm y\bigr)^2 ,
\]
जहाँ बीच का चरण कोशी--श्वार्ट्ज़ है; और समानता $\langle x, y\rangle = \norm
x\norm y$ बाध्य कर देती है, अर्थात् \emph{धनात्मक} समानता वाली स्थिति, यानी
ऋणेतर अनुपात के साथ अनुपातिता --- ज्यामितीय रूप से त्रिभुज तभी अपभ्रष्ट होता
है जब दोनों सदिश एक ही दिशा में हों। अंतिम दोनों सर्वसमिकाएँ सीधे प्रसार हैं
(दूसरी, अर्थात् \emph{ध्रुवण सर्वसमिका}, मानक से गुणन को फिर से प्राप्त कर
लेती है)।
\end{proof}

\section{लांबिकता}

\begin{definition}\label{def:b1:euclid:orthogonal}
$x \perp y$ तब जब $\langle x, y \rangle = 0$। कोई कुल \emph{लांबिक} कहलाता
है जब उसके सदिश युग्मशः लांबिक हों, और
\emph{लांबिक-प्रसामान्य}\index{लांबिक-प्रसामान्य कुल} जब इसके अतिरिक्त हर एक
का मानक $1$ हो। किसी \omterm{def:b1:vspaces:subspace}{उपसमष्टि} $F$ का \emph{लांबिक \omterm{def:b1:vspaces:sum}{पूरक}} है
\[
F^{\perp} = \{x \in E : \forall y \in F,\ \langle x, y\rangle =
0\},
\]
जो $E$ की एक \omterm{def:b1:vspaces:subspace}{उपसमष्टि} है।\index{लांबिक पूरक}
\end{definition}

\begin{proposition}\label{prop:b1:euclid:orthfree}
(पाइथागोरस\index{पाइथागोरस प्रमेय}) यदि $x \perp y$, तो $\norm{x+y}^2 =
\norm x^2 + \norm y^2$। अशून्य सदिशों का कोई \omterm{def:b1:euclid:orthogonal}{लांबिक} कुल \omterm{def:b1:vspaces:free}{स्वतंत्र} होता है।
किसी \omterm{def:b1:euclid:orthogonal}{लांबिक-प्रसामान्य} \omterm{def:b1:vspaces:free}{आधार} $(e_1, \dots, e_n)$ में \omterm{prop:b1:vspaces:coordinates}{निर्देशांक} और गुणन ये
हैं
\[
x = \sum_{i} \langle x, e_i\rangle\, e_i,
\qquad
\langle x, y \rangle = \sum_i \langle x, e_i\rangle \langle y,
e_i\rangle,
\qquad
\norm x^2 = \sum_i \langle x, e_i\rangle^2 .
\]
\end{proposition}

\begin{proof}
पाइथागोरस: प्रसार कीजिए। स्वतंत्रता: किसी शून्य संयोजन का $\langle\,\cdot\,,
x_j\rangle$ लीजिए: $\lambda_j \norm{x_j}^2 = 0$। \omterm{prop:b1:vspaces:coordinates}{निर्देशांक}: $x = \sum
\lambda_i e_i$ लिखिए और $e_j$ के साथ गुणन लीजिए: $\lambda_j = \langle x,
e_j\rangle$; और दोनों सूत्र द्विरैखिकता से आ जाते हैं।
\end{proof}

\begin{example}[लांबिक-प्रसामान्य निर्देशांक, पार्सेवाल जाँच के साथ]
\label{ex:b1:euclid:parsevalcheck}
\cref{exo:b1:euclid:3} के \omterm{def:b1:euclid:orthogonal}{लांबिक-प्रसामान्य} \omterm{def:b1:vspaces:free}{आधार} में $x = (1, 2, 3)$ का
प्रसार कीजिए,
\[
e_1 = \tfrac{1}{\sqrt2}(1,1,0), \quad
e_2 = \tfrac{1}{\sqrt6}(1,-1,2), \quad
e_3 = \tfrac{1}{\sqrt3}(-1,1,1).
\]
कोई निकाय हल नहीं करना पड़ा --- केवल तीन \omterm{def:b1:euclid:def}{अंतर्गुणन}:
\[
\langle x, e_1\rangle = \frac{3}{\sqrt2},
\qquad
\langle x, e_2\rangle = \frac{1 - 2 + 6}{\sqrt6} =
\frac{5}{\sqrt6},
\qquad
\langle x, e_3\rangle = \frac{-1 + 2 + 3}{\sqrt3} =
\frac{4}{\sqrt3}.
\]
प्रतिज्ञप्ति के मानक-सूत्र से प्रमाणन:
\[
\frac{9}{2} + \frac{25}{6} + \frac{16}{3}
= \frac{27 + 25 + 32}{6} = 14 = \norm{x}^2 = 1 + 4 + 9 .
\]
निर्देशांकों के वर्गों के योग वाली यह जाँच (एक परिमित पार्सेवाल सर्वसमिका)
कुछ सेकंड लेती है और चिह्न तथा प्रसामान्यन की त्रुटियाँ लगभग निश्चित रूप से
पकड़ लेती है --- जब भी कोई \omterm{def:b1:euclid:orthogonal}{लांबिक-प्रसामान्य} प्रसार संगणित करें, इसे आदत बना
लीजिए; और \cref{ex:b1:euclid:trigortho} के फूरिये गुणांकों के लिए उसका
अनंत-विमीय रूप स्नातक वर्ष~3 के खंड की एक प्रमेय है।
\end{example}

\begin{theorem}[ग्राम--श्मिट]\label{thm:b1:euclid:gramschmidt}
हर \omterm{def:b1:euclid:def}{यूक्लिडीय समष्टि} में \omterm{def:b1:euclid:orthogonal}{लांबिक-प्रसामान्य} \omterm{def:b1:vspaces:free}{आधार} होते हैं। स्पष्ट रूप से, किसी
भी \omterm{def:b1:vspaces:free}{आधार} $(v_1, \dots, v_n)$ से आरंभ करके यह नुस्ख़ा\index{ग्राम--श्मिट
प्रक्रिया}
\[
w_k = v_k - \sum_{i=1}^{k-1} \langle v_k, e_i\rangle\, e_i ,
\qquad
e_k = \frac{w_k}{\norm{w_k}}
\]
एक \omterm{def:b1:euclid:orthogonal}{लांबिक-प्रसामान्य} \omterm{def:b1:vspaces:free}{आधार} $(e_1, \dots, e_n)$ बना देता है, जिसमें हर $k$ के
लिए $\operatorname{Vect}(e_1, \dots, e_k) = \operatorname{Vect}(v_1, \dots,
v_k)$।
\end{theorem}

\begin{proof}
$k$ पर आगमन। मान लीजिए $(e_1, \dots, e_{k-1})$ \omterm{def:b1:euclid:orthogonal}{लांबिक-प्रसामान्य} हैं और
$\operatorname{Vect}(v_1, \dots, v_{k-1})$ जनित करते हैं: तब सदिश $w_k$ रचना
से ही हर $e_j$ ($j < k$) के \omterm{def:b1:euclid:orthogonal}{लांबिक} है ($\langle w_k, e_j \rangle = \langle
v_k, e_j\rangle - \langle v_k, e_j\rangle$), और $w_k \neq 0$, क्योंकि $v_k
\notin \operatorname{Vect}(v_1, \dots, v_{k-1})$। प्रसामान्यन लांबिकता बनाए
रखता है; और \omterm{def:b1:vspaces:span}{जनित समष्टि} वाला \omterm{def:b1:logic:statement}{कथन} इसलिए सत्य है कि $e_k$ $v_k$ तथा पहले के
$e_i$ का संयोजन है, वह भी व्युत्क्रमणीय रूप से।
\end{proof}

\begin{example}[बहुपदों पर ग्राम--श्मिट, पूरा]
\label{ex:b1:euclid:legendre}
$\langle P, Q\rangle = \int_0^1 PQ$ के साथ $\R_2[X]$ में $(1, X, X^2)$ का
लांबिक-प्रसामान्यन कीजिए। \emph{चरण 1}: $\norm{1}^2 = 1$, अतः $e_1 = 1$।
\emph{चरण 2}: $w_2 = X - \langle X, 1\rangle\,1 = X - \frac12$, और
$\norm{w_2}^2 = \int_0^1\bigl(x - \frac12\bigr)^2 \dd x = \frac1{12}$: $e_2
= \sqrt{12}\,\bigl(X - \frac12\bigr)$। \emph{चरण 3}: $\langle X^2,
e_1\rangle = \frac13$ और
\[
\langle X^2, e_2\rangle
= \sqrt{12}\int_0^1 x^2\Bigl(x - \frac12\Bigr)\dd x
= \frac{\sqrt{12}}{12},
\qquad\text{अतः}\qquad
w_3 = X^2 - \frac13 - \Bigl(X - \frac12\Bigr)
= X^2 - X + \frac16 .
\]
उसका मानक \cref{exo:b1:euclid:9} में संगणित हुआ था: $\norm{w_3}^2 =
\frac1{180}$, जिससे $e_3 = \sqrt{180}\,\bigl(X^2 - X + \frac16\bigr)$। \omterm{def:b1:poly:def}{बहुपद}
$1$, $X - \frac12$, $X^2 - X + \frac16$ मापन तक \omterm{prop:b1:reals:intervals}{अंतराल} $\intcc{0}{1}$ के
पहले \emph{लजांद्र \omterm{def:b1:poly:def}{बहुपद}} हैं; और रचना एक-एक घात करके आगे चलती जाती है, जहाँ
हर नया \omterm{def:b1:poly:def}{बहुपद} अपने सभी पूर्ववर्तियों के \omterm{def:b1:euclid:orthogonal}{लांबिक} होता है। ध्यान दीजिए कि
कलनविधि पहले का काम फिर से काम में ले लेती है: चरण 3 पर घटाया गया \omterm{def:b1:linmaps:projection}{प्रक्षेप}
ठीक वही $X^2$ का सर्वोत्तम एफ़ीन सन्निकटन है जो
\cref{ex:b1:euclid:bestapprox} में मिला था --- ग्राम--श्मिट पुनरावृत्त
\omterm{thm:b1:euclid:projection}{लांबिक प्रक्षेप} \emph{ही} है।
\end{example}

\begin{theorem}[लांबिक प्रक्षेप]\label{thm:b1:euclid:projection}
मान लीजिए $F$ \omterm{def:b1:euclid:def}{यूक्लिडीय समष्टि} $E$ की \omterm{def:b1:vspaces:subspace}{उपसमष्टि} है। तब
\[
E = F \oplus F^{\perp},
\]
और $F$ पर संबद्ध \omterm{def:b1:linmaps:projection}{प्रक्षेप} $p_F$ (अर्थात् \emph{लांबिक प्रक्षेप}\index{लांबिक
प्रक्षेप}) $F$ के किसी भी \omterm{def:b1:euclid:orthogonal}{लांबिक-प्रसामान्य} \omterm{def:b1:vspaces:free}{आधार} $(e_1, \dots, e_k)$ में
$p_F(x) = \sum_i \langle x, e_i\rangle e_i$ से मिलता है। वही $F$ तक की दूरी
को साकार करता है: सभी $y \in F$ के लिए,
\[
\norm{x - p_F(x)} \leq \norm{x - y},
\]
और समानता केवल $y = p_F(x)$ के लिए; इसे $d(x, F) = \norm{x - p_F(x)}$ लिखा
जाता है।
\end{theorem}

\begin{proof}
$F$ का कोई \omterm{def:b1:euclid:orthogonal}{लांबिक-प्रसामान्य} \omterm{def:b1:vspaces:free}{आधार} $(e_i)_{i \leq k}$ लीजिए ($F$ के भीतर
\cref{thm:b1:euclid:gramschmidt}) और $\pi(x) = \sum \langle x, e_i\rangle
e_i \in F$ रखिए। तब हर $j$ के लिए $x - \pi(x) \perp e_j$ (वही निरसन जैसा ऊपर
था), अतः $x - \pi(x) \in F^\perp$: $E = F + F^\perp$। और $F \cap F^\perp =
\{0\}$: ऐसा सदिश $\langle x, x\rangle = 0$ संतुष्ट करता है। अतः योग
प्रत्यक्ष है और $\pi = p_F$।

दूरी: $y \in F$ के लिए $x - y = (x - p_F(x)) + (p_F(x) - y)$ का विघटन कीजिए,
जिसके टुकड़े \omterm{def:b1:euclid:orthogonal}{लांबिक} हैं ($F^\perp$ और $F$); पाइथागोरस:
\[
\norm{x - y}^2 = \norm{x - p_F(x)}^2 + \norm{p_F(x) - y}^2
\geq \norm{x - p_F(x)}^2,
\]
और समानता तभी जब $y = p_F(x)$।
\end{proof}

\begin{example}[प्रक्षेप कभी लंबा नहीं करते]
\label{ex:b1:euclid:bessel}
\omterm{thm:b1:logic:partition}{विभाजन} $x = p_F(x) + (x - p_F(x))$ पर पाइथागोरस लगाने से:
\[
\norm{p_F(x)}^2 = \norm x^2 - \norm{x - p_F(x)}^2 \leq
\norm x^2 ,
\]
और समानता तभी जब $x \in F$। $F$ के किसी \omterm{def:b1:euclid:orthogonal}{लांबिक-प्रसामान्य} \omterm{def:b1:vspaces:free}{आधार} $(e_1, \dots,
e_k)$ में यह $\sum_{i \leq k}\langle x, e_i\rangle^2 \leq \norm x^2$ कहता है
(एक \emph{बेसेल असमिका}): कितनी भी \omterm{def:b1:euclid:orthogonal}{लांबिक-प्रसामान्य} दिशाएँ नाप लीजिए,
निर्देशांकों के वर्ग कभी लंबाई के वर्ग से अधिक नहीं होते --- और जब वह कुल
पूरा \omterm{def:b1:vspaces:free}{आधार} हो, तब \cref{ex:b1:euclid:parsevalcheck} की यथार्थ समानता से तुलना
कीजिए। यही एक-पंक्ति की असमिका स्नातक वर्ष~3 के खंड में फूरिये गुणांकों को
योग्य बनाती है; और यहीं वह यह भी समझा देती है कि किसी न्यूनतम-वर्ग आरोपण में
और आधार-फलन जोड़ने से अवशिष्ट केवल घट ही सकता है।
\end{example}

\begin{example}[सर्वोत्तम वर्ग-माध्य सन्निकटन]\label{ex:b1:euclid:bestapprox}
$\langle f, g\rangle = \int_0^1 fg$ के साथ $C(\intcc{0}{1})$ में, संबद्ध
(वर्ग-माध्य) दूरी में $f(x) = x^2$ के सबसे निकट $\leq 1$ घात का \omterm{def:b1:poly:def}{बहुपद}
$p_F(f)$ है, जहाँ $F = \R_1[X]$। $(1, X)$ पर ग्राम--श्मिट: $e_1 = 1$, $w_2 =
X - \frac12$, $\norm{w_2}^2 = \int_0^1 (x - \frac12)^2 = \frac{1}{12}$, $e_2
= \sqrt{12}\,(X - \tfrac12)$। तब
\[
p_F(f) = \langle f, e_1\rangle e_1 + \langle f, e_2\rangle e_2
= \frac13 + 12\Bigl(\int_0^1 x^2\bigl(x - \tfrac12\bigr)\dd
x\Bigr)\bigl(X - \tfrac12\bigr)
= X - \frac{1}{6},
\]
$\int_0^1 x^2(x - \frac12)\dd x = \frac14 - \frac16 = \frac{1}{12}$ का
प्रयोग करते हुए। अर्थात् ``न्यूनतम वर्ग'' का विचार रैखिक बीजगणित की एक
पंक्ति में।
\end{example}

\begin{method}[किसी दूरी $d(x, F)$ तक के तीन रास्ते]
\label{met:b1:euclid:distance}
\begin{enumerate}
  \item \emph{$F$ का \omterm{def:b1:euclid:orthogonal}{लांबिक-प्रसामान्य} \omterm{def:b1:vspaces:free}{आधार}}: तब $p_F(x) = \sum_i \langle x,
  e_i\rangle e_i$, और पाइथागोरस से,
        \[
        d(x, F)^2 = \norm{x}^2 - \norm{p_F(x)}^2
        = \norm x^2 - \sum_i \langle x, e_i\rangle^2 ,
        \]
        जो प्रायः स्वयं $x - p_F(x)$ संगणित करने से सस्ता पड़ता है।
  \item \emph{प्रसामान्य समीकरण}: $F$ के किसी भी जनक कुल के साथ $p$ के
  गुणांकों के लिए $\langle x - p, v_j\rangle = 0$ हल कीजिए
  (\cref{exo:b1:euclid:5}) --- किसी लांबिक-प्रसामान्यन की आवश्यकता नहीं।
  \item \emph{\omterm{def:b1:vspaces:sum}{पूरक} के रास्ते}: यदि $F^\perp$ $F$ से छोटा है (जैसे $F$ कोई
  \omterm{def:b1:linmaps:forms}{अधिसमतल} और $F^\perp$ कोई रेखा $\operatorname{Vect}(n)$), तो उसके बदले
  $F^\perp$ पर \omterm{def:b1:linmaps:projection}{प्रक्षेप} कीजिए:
        \[
        d(x, F) = \norm{p_{F^\perp}(x)} =
        \frac{\abs{\langle x, n\rangle}}{\norm n} ,
        \]
        यही समतल-तक-की-दूरी वाला शास्त्रीय सूत्र है
        (\cref{exo:b1:multivar:8} उसका प्रयोग करता है)।
\end{enumerate}
रास्ता 3 एक व्यापक प्रतिवर्त की विशेष स्थिति है: $F$ और $F^\perp$ में से
जिसकी विमा छोटी हो, सदा उसी पर \omterm{def:b1:linmaps:projection}{प्रक्षेप} कीजिए।
\end{method}

\begin{example}[त्रिकोणमितीय लांबिकता: फूरिये की एक झलक]
\label{ex:b1:euclid:trigortho}
$\langle f, g\rangle = \frac1\pi\int_0^{2\pi} fg$ के साथ
$C(\intcc{0}{2\pi})$ पर, कुल
\[
\Bigl(\frac{1}{\sqrt2},\ \cos x,\ \sin x,\ \cos 2x,\ \sin 2x,\
\dots\Bigr)
\]
\omterm{def:b1:euclid:orthogonal}{लांबिक-प्रसामान्य} है: उदाहरण के लिए $p \neq q$ के लिए $\langle \cos px, \cos
qx\rangle = \frac1\pi\int_0^{2\pi}\cos px\cos qx\,\dd x = 0$ (गुणनफल को
$\frac12[\cos(p{-}q)x + \cos(p{+}q)x]$ में \omterm{def:b1:linmaps:def}{रैखिक} कीजिए और पूरे आवर्तों पर
समाकलन कीजिए), जबकि $\frac1\pi\int_0^{2\pi}\cos^2 px\,\dd x = 1$। अतः इन
फलनों में से पहले $2N + 1$ की \omterm{def:b1:vspaces:span}{जनित समष्टि} पर \omterm{thm:b1:euclid:projection}{लांबिक प्रक्षेप} के \omterm{prop:b1:vspaces:coordinates}{निर्देशांक}
$\langle f, e_i\rangle$ होते हैं --- अर्थात् कोज्याओं और ज्याओं के सापेक्ष
\omterm{thm:b1:integration:def}{समाकल}। ये $f$ के \emph{फूरिये गुणांक} हैं, और वह \omterm{def:b1:linmaps:projection}{प्रक्षेप} उसका सर्वोत्तम
वर्ग-माध्य त्रिकोणमितीय सन्निकटन है; उनके अभिसरण का अध्ययन स्नातक वर्ष~3 का
खंड करता है। सारा काम लांबिकता कर रही है: गुणांकों के सूत्र अक्षरशः
\cref{thm:b1:euclid:projection} ही हैं।
\end{example}

\section{समतल के समदूरक}

\begin{definition}\label{def:b1:euclid:isometry}
किसी \omterm{def:b1:euclid:def}{यूक्लिडीय समष्टि} की अंतःसमाकारिता $u$ \emph{समदूरक}\index{समदूरक} (अथवा
\omterm{def:b1:euclid:orthogonal}{लांबिक} \omterm{def:b1:logic:map}{प्रतिचित्रण}) तब कहलाती है जब वह मानक सुरक्षित रखे: सभी $x$ के लिए
$\norm{u(x)} = \norm x$ --- तुल्य रूप से (ध्रुवण से) वह \omterm{def:b1:euclid:def}{अंतर्गुणन} सुरक्षित
रखती है; और तुल्य रूप से किसी \omterm{def:b1:euclid:orthogonal}{लांबिक-प्रसामान्य} \omterm{def:b1:vspaces:free}{आधार} में उसका आव्यूह $A$
$A^{\mathsf T} A = I$ संतुष्ट करता है। समदूरक एक \omterm{def:b1:structures:group}{समूह} बनाते हैं, अर्थात्
\omterm{def:b1:euclid:orthogonal}{लांबिक} \omterm{def:b1:structures:group}{समूह} $O(E)$।
\end{definition}

\begin{example}[समदूरक को देखते ही पहचानना]
\label{ex:b1:euclid:recognize}
क्या $A = \dfrac15\begin{pmatrix} 3 & -4\\ 4 & 3\end{pmatrix}$ \omterm{def:b1:euclid:orthogonal}{लांबिक} है?
स्तंभ: मानक $\frac15\sqrt{9 + 16} = 1$ और $\frac15\sqrt{16 + 9} = 1$; गुणन
$\frac1{25}(3\cdot(-4) + 4\cdot3) = 0$। हाँ --- और $\det A = \frac{9 +
16}{25} = 1$, अतः वह $\cos\theta = \frac35$, $\sin\theta = \frac45$ वाला
घूर्णन $R_\theta$ है (``$3$-$4$-$5$ घूर्णन'', जिसका कोण $\pi$ का कोई
उल्लेखनीय भिन्नांश नहीं है)। इसके विपरीत $B =
\frac{1}{\sqrt2}\begin{pmatrix} 1 & 1\\ 0 & 1\end{pmatrix}$ देखने में तो
इकाई-सारणिक जैसा लगता है, पर उसका पहला स्तंभ इकाई नहीं है
($\frac1{\sqrt2}$): अतः \omterm{def:b1:euclid:orthogonal}{लांबिक} नहीं --- अकेला सारणिक $\pm1$ कुछ भी प्रमाणित
नहीं करता, स्तंभ जाँचने ही पड़ते हैं।
\end{example}

\begin{theorem}[समतल के समदूरक]\label{thm:b1:euclid:plane}
किसी यूक्लिडीय समतल के \omterm{def:b1:euclid:orthogonal}{लांबिक-प्रसामान्य} \omterm{def:b1:vspaces:free}{आधार} में समदूरकों के आव्यूह ठीक ये
हैं
\[
R_\theta = \begin{pmatrix}
\cos\theta & -\sin\theta\\ \sin\theta & \cos\theta
\end{pmatrix}
\quad (\text{घूर्णन, कोण } \theta,\ \det = 1),
\]
\[
S_\theta = \begin{pmatrix}
\cos\theta & \sin\theta\\ \sin\theta & -\cos\theta
\end{pmatrix}
\quad (\det = -1),
\]
जिनमें दूसरा उस रेखा में परावर्तन है जो पहले आधार-सदिश के साथ कोण
$\frac\theta2$ बनाती है।
\end{theorem}

\begin{proof}
मान लीजिए $A^{\mathsf T}A = I$ के साथ $A = \begin{pmatrix} a & c\\ b &
d\end{pmatrix}$: स्तंभ इकाई और \omterm{def:b1:euclid:orthogonal}{लांबिक} हैं। पहला स्तंभ किसी $\theta$ के लिए
$(\cos\theta, \sin\theta)$ है; और दूसरा, जो इकाई तथा उसके \omterm{def:b1:euclid:orthogonal}{लांबिक} है,
$\pm(-\sin\theta, \cos\theta)$ है। चिह्न $+$ $R_\theta$ देता है; चिह्न $-$
$S_\theta$ देता है। जाँच लीजिए कि $S_\theta^2 = I$, और यह कि सदिश
$(\cos\frac\theta2, \sin\frac\theta2)$ अचल रहता है जबकि उसका \omterm{def:b1:euclid:orthogonal}{लांबिक} उलट जाता
है: अर्थात् एक परावर्तन। (और $R_\alpha R_\beta = R_{\alpha+\beta}$:
घूर्णन-समूह ही कोण-समूह है --- \cref{thm:b1:complex:funceq} से तुलना कीजिए।)
\end{proof}

\begin{omfigure}
\begin{tikzpicture}[scale=1.15]
  \draw[->, thin] (-1.2,0) -- (2.9,0) node[below] {$x$};
  \draw[->, thin] (0,-1.2) -- (0,2.6) node[left] {$y$};
  \draw[omDef, very thick] (-0.9,0) -- (2.6,0)
    node[below, pos=0.9, font=\small] {$r_1$ का अक्ष};
  \draw[omThm, very thick] (-0.9,-0.9) -- (2.3,2.3)
    node[right, pos=0.97, font=\small] {$r_2$ का अक्ष};
  \fill (2,0.5) circle (0.05) node[right, font=\small] {$M$};
  \fill (2,-0.5) circle (0.05)
    node[right, font=\small] {$r_1(M)$};
  \fill (-0.5,2) circle (0.05)
    node[above, font=\small] {$r_2(r_1(M))$};
  \draw[omDef, dashed] (2,0.5) -- (2,-0.5);
  \draw[omThm, dashed] (2,-0.5) -- (-0.5,2);
  \draw[omProp, thick, ->] (0.9,0) arc (0:45:0.9);
  \node[omProp, font=\small] at (1.15,0.32) {$\frac\pi4$};
\end{tikzpicture}

{\small दो परावर्तन मिलकर घूर्णन बनाते हैं: $M = (2, 0.5)$ को $x$-अक्ष में
परावर्तित कीजिए, फिर रेखा $y = x$ में, तो वह $(-0.5, 2)$ पर पहुँचता है ---
अर्थात् मूल बिंदु के परितः $\frac\pi2$ कोण के घूर्णन के अंतर्गत $M$ का
\omterm{def:b1:linmaps:kerim}{प्रतिबिंब}, और यह कोण दोनों अक्षों के बीच के कोण $\frac\pi4$ का \emph{दुगुना}
है। सप्ताहांत समस्या इसी चित्र को समतल के सभी समदूरकों के संयोजन-नियम में
बदल देती है।}
\end{omfigure}

\begin{remark}[सामान्य भूलें]
\label{rem:b1:euclid:pitfalls}
\emph{प्रक्षेप-सूत्र को \omterm{def:b1:euclid:orthogonal}{लांबिक-प्रसामान्य} \omterm{def:b1:vspaces:free}{आधार} चाहिए}: $F$ के केवल जनक कुल
$(v_i)$ के लिए योग $\sum_i\langle x, v_i\rangle v_i$ $p_F(x)$ \emph{नहीं} है
($F = \R^2$, $v_1 = e_1$, $v_2 = e_1 + e_2$ पर परखिए); और \omterm{def:b1:euclid:orthogonal}{लांबिक-प्रसामान्य}
न होने वाले कुल के साथ उसके बदले प्रसामान्य समीकरण हल कीजिए
(\cref{met:b1:euclid:distance}~(2))। \emph{\omterm{def:b1:vspaces:free}{स्वतंत्र} होने के लिए \omterm{def:b1:euclid:orthogonal}{लांबिक} कुलों
को $0$ से बचना चाहिए}: शून्य सदिश हर चीज़ के \omterm{def:b1:euclid:orthogonal}{लांबिक} है, स्वयं अपने भी ---
अतः \cref{prop:b1:euclid:orthfree} में स्वतंत्रता के लिए सदिशों का अशून्य
होना आवश्यक है। \emph{$F^\perp$ \omterm{def:b1:euclid:def}{अंतर्गुणन} पर निर्भर करता है}: $\R_1[X]$ में
$\int_0^1 PQ$ के लिए $\operatorname{Vect}(X)$ का \omterm{def:b1:vspaces:sum}{पूरक} अचर फलन नहीं, बल्कि
$\operatorname{Vect}(1 - \frac32 X)$ है --- $\int_0^1 x(1 - \frac32 x) =
\frac12 - \frac12 = 0$ संगणित कीजिए; और जब तक गुणन का नाम न लिया जाए,
``लंबवत'' निरर्थक है। \emph{$\norm{x + y}$ का \omterm{def:b1:linmaps:def}{रैखिक} प्रसार मत कीजिए}: सही
सर्वसमिका $\norm{x+y}^2 = \norm x^2 + 2\langle x, y\rangle + \norm y^2$ है;
बीच वाला पद केवल लांबिकता के अंतर्गत शून्य होता है (पाइथागोरस), और त्रिभुज
असमिका एक \emph{असमिका} है। \emph{इकाई सदिशों को इकाई सदिशों पर भेजना
पर्याप्त नहीं है}: $u(x, y) = (x + y,\ 0)$ दोनों विहित आधार-सदिशों को इकाई
सदिश $(1, 0)$ पर भेजता है, फिर भी $\norm{u(1,1)} = 2 \neq \sqrt2$: अतः कोई
\omterm{def:b1:euclid:isometry}{समदूरक} नहीं। परिभाषा \emph{सभी} $x$ के लिए $\norm{u(x)} = \norm x$ माँगती
है; आव्यूह की भाषा में $A^{\mathsf T}A = I$, अर्थात् ऐसे स्तंभ जो इकाई
\emph{और युग्मशः \omterm{def:b1:euclid:orthogonal}{लांबिक}} हों --- दोनों शर्तें, एक साथ जाँची हुई।
\end{remark}

\begin{remark}[अंतर्गुणन कहाँ जाता है]
\label{rem:b1:euclid:whereused}
\omterm{thm:b1:euclid:projection}{लांबिक प्रक्षेप} इस अध्याय की सबसे अधिक प्रयुक्त प्रमेय है: वही न्यूनतम
वर्गों के नीचे है (\cref{ch:b1:multivar} की सप्ताहांत समस्या उसी पर
समाश्रयण-रेखाएँ बनाती है), फूरिये गुणांकों के नीचे
(\cref{ex:b1:euclid:trigortho}), और \cref{exo:b1:euclid:5} के प्रसामान्य
समीकरणों के नीचे, जिन्हें संख्यात्मक विश्लेषण बड़े पैमाने पर हल करता है।
वर्गीकरण $R_\theta / S_\theta$ नीचे पूरा किया गया है: सप्ताहांत समस्या समतल
के \emph{सभी} दूरी-संरक्षी रूपांतरणों का वर्गीकरण करती है, चाहे वे \omterm{def:b1:linmaps:def}{रैखिक} हों
या न हों, और उनके परिमित समूहों का भी --- अर्थात् पुष्पिकाओं और सम बहुभुजों
का गणित। स्नातक वर्ष~2 के खंड में \omterm{def:b1:euclid:def}{अंतर्गुणन} अभिलक्षणिक-मान सिद्धांत से मिलता
है (सममित आव्यूह, द्विघातीय रूप); और स्नातक वर्ष~3 में अनंत-विमीय यूक्लिडीय
ज्यामिति हिल्बर्ट समष्टि का सिद्धांत बन जाती है।
\end{remark}

\begin{remark}[पुस्तक 3 के भीतर के परिदृश्य]
\label{rem:b1:euclid:perspectives}
इस अध्याय से दो पुल निकलते हैं। \emph{पीछे की ओर}, रैखिक बीजगणित तक:
\cref{exo:b1:euclid:11} का ग्राम आव्यूह अंतर्गुणनों को \cref{ch:b1:det} की
सारणिक-मशीनरी में बाँध देता है, और \omterm{thm:b1:euclid:projection}{लांबिक प्रक्षेप} \cref{ch:b1:linmaps} का
वही विशेष प्रक्षेपक है जिसकी अष्टि $F^\perp$ है --- उसका सारा बीजगणित ($p^2
= p$, $s = 2p - \mathrm{id}$) अक्षरशः लागू होता है, और अब बोनस यह है कि
$\norm{x - p(x)}$ एक \emph{दूरी} है। \emph{आगे की ओर}, विश्लेषण तक:
\cref{ch:b1:curves} चाप की लंबाई इसी अध्याय के मानक से नापता है और उसके
समदूरकों के बिना कुछ भी वर्गीकृत नहीं करता; \cref{ch:b1:multivar} प्रवणता को
कोशी--श्वार्ट्ज़ के द्वारा पढ़ता है (तीव्रतम आरोहण) और खंड को न्यूनतम वर्गों
के साथ बंद करता है, जो $\R^n$ के किसी आँकड़ा-सदिश पर लगाया गया
\cref{thm:b1:euclid:projection} ही है। \omterm{def:b1:euclid:def}{अंतर्गुणन} वही बिंदु है जहाँ इस पुस्तक
का बीजगणित और उसका विश्लेषण अंततः मिलते हैं।
\end{remark}

\section{अभ्यास}

\begin{exercise}[$\star$]\label{exo:b1:euclid:1}
विहित $\R^3$ में: $\langle u, v\rangle$, $\norm u$, $\norm v$ तथा $u = (1,
2, 2)$ और $v = (2, -2, 1)$ के बीच का कोण संगणित कीजिए। कोशी--श्वार्ट्ज़
संख्यात्मक रूप से सत्यापित कीजिए।
\end{exercise}

\begin{exercise}[$\star$]\label{exo:b1:euclid:2}
किसी भी \omterm{def:b1:euclid:def}{यूक्लिडीय समष्टि} में समांतर चतुर्भुज सर्वसमिका $\norm{x+y}^2 +
\norm{x-y}^2 = 2\norm x^2 + 2\norm y^2$ सिद्ध कीजिए, और उसका प्रयोग करके
दिखाइए कि $\R^2$ पर \omterm{def:b1:reals:bounds}{उच्चतम} मानक $\norm{(x,y)}_\infty = \max(\abs x, \abs y)$
किसी \omterm{def:b1:euclid:def}{अंतर्गुणन} से नहीं आता।
\end{exercise}

\begin{exercise}[$\star$]\label{exo:b1:euclid:3}
विहित $\R^3$ में $\bigl((1,1,0), (1,0,1), (0,1,1)\bigr)$ पर ग्राम--श्मिट
लगाइए।
\end{exercise}

\begin{exercise}[$\star$]\label{exo:b1:euclid:4}
$\R^3$ में मान लीजिए $F = \operatorname{Vect}\bigl((1,1,1)\bigr)$। $F^\perp$
निर्धारित कीजिए (समीकरण और \omterm{def:b1:vspaces:free}{आधार}), विहित \omterm{def:b1:vspaces:free}{आधार} में $p_F$ का आव्यूह, तथा
$d\bigl((1, 2, 3), F\bigr)$।
\end{exercise}

\begin{exercise}[$\star\star$]\label{exo:b1:euclid:5}
(प्रसामान्य समीकरण) मान लीजिए $F =
\operatorname{Vect}\bigl((1,0,1),(0,1,1)\bigr) \subseteq \R^3$ और $x = (1,
1, 4)$। दोनों जनकों के लिए $\langle x - p, v \rangle = 0$ हल करके $p_F(x)$
संगणित कीजिए ($p = \alpha(1,0,1) + \beta(0,1,1)$), फिर $d(x, F)$। यहाँ
ग्राम--श्मिट अनावश्यक क्यों है?
\end{exercise}

\begin{exercise}[$\star\star$]\label{exo:b1:euclid:6}
$\intcc{0}{1}$ पर \omterm{def:b1:continuity:continuous}{संतत} $f$ के लिए सिद्ध कीजिए
\[
\Bigl(\int_0^1 f\Bigr)^{\!2} \leq \int_0^1 f^2 ,
\]
समानता की स्थिति के साथ, $C(\intcc{0}{1})$ में कोशी--श्वार्ट्ज़ के एक उदाहरण
के रूप में। फिर वास्तविक $a_i$ के लिए $\bigl(\sum_{i=1}^n a_i\bigr)^2 \leq n
\sum a_i^2$ सिद्ध कीजिए।
\end{exercise}

\begin{exercise}[$\star\star$]\label{exo:b1:euclid:7}
सिद्ध कीजिए कि किसी \omterm{def:b1:euclid:def}{यूक्लिडीय समष्टि} की हर \omterm{def:b1:vspaces:subspace}{उपसमष्टि} $F$ के लिए:
$(F^{\perp})^{\perp} = F$ और $\dim F^{\perp} = \dim E - \dim F$।
\end{exercise}

\begin{exercise}[$\star\star$]\label{exo:b1:euclid:8}
निम्नलिखित आव्यूहों वाले \omterm{def:b1:euclid:isometry}{समतल-समदूरक} पहचानिए
\[
A = \frac{1}{\sqrt 2}\begin{pmatrix} 1 & -1\\ 1 & 1\end{pmatrix},
\qquad
B = \frac{1}{5}\begin{pmatrix} 3 & 4\\ 4 & -3\end{pmatrix}
\]
(प्रकार, कोण अथवा अक्ष)। $A^8$ और $B^2$ बिना आव्यूह गुणा किए संगणित कीजिए।
\end{exercise}

\begin{exercise}[$\star\star\star$]\label{exo:b1:euclid:9}
(\omterm{def:b1:linmaps:projection}{प्रक्षेप} के रूप में न्यूनतम) संगणित कीजिए
\[
\min_{(a, b) \in \R^2} \int_0^1 \bigl(x^2 - a - bx\bigr)^2 \dd x ,
\]
\cref{ex:b1:euclid:bestapprox} का प्रयोग करते हुए: न्यूनतम $\norm{f -
p_F(f)}^2$ है, और वह $f = X^2$, $F = \R_1[X]$ के लिए।
\end{exercise}

\begin{exercise}[$\star\star\star$]\label{exo:b1:euclid:10}
मान लीजिए $u$ \omterm{def:b1:euclid:def}{यूक्लिडीय समष्टि} $E$ का कोई \omterm{def:b1:euclid:isometry}{समदूरक} है। सिद्ध कीजिए कि $\ker(u
- \mathrm{id}) \perp \operatorname{im}(u - \mathrm{id})$, और $E = \ker(u -
\mathrm{id}) \oplus \operatorname{im}(u - \mathrm{id})$ निकालिए। \emph{(गुणन
के संरक्षण का प्रयोग करते हुए $u(y) = y$ के लिए $\langle x - u(x), y\rangle$
संगणित कीजिए।)}
\end{exercise}

\begin{exercise}[$\star\star$]\label{exo:b1:euclid:11}
(ग्राम आव्यूह) किसी \omterm{def:b1:euclid:def}{यूक्लिडीय समष्टि} के सदिशों $v_1, \dots, v_k$ के लिए मान
लीजिए $G = \bigl(\langle v_i, v_j\rangle\bigr)_{1 \leq i, j \leq k}$ उनका
\emph{ग्राम आव्यूह} है।
\begin{enumerate}
  \item सिद्ध कीजिए कि $(v_1, \dots, v_k)$ \omterm{def:b1:vspaces:free}{स्वतंत्र} है तभी जब $G$
  व्युत्क्रमणीय हो। \emph{(यदि $Gc = 0$, तो $\norm{\sum_i c_i v_i}^2$ संगणित
  कीजिए।)}
  \item $\langle P, Q\rangle = \int_0^1 PQ$ के साथ $\R_2[X]$ में $(1, X,
  X^2)$ का ग्राम आव्यूह संगणित कीजिए, \cref{ch:b1:det} की सप्ताहांत समस्या
  वाला \omterm{pb:b1:det:1}{हिल्बर्ट आव्यूह} $H_3$ पहचानिए, और $\det H_3 = \frac1{2160} \neq 0$ से
  स्वतंत्रता का निष्कर्ष निकालिए।
\end{enumerate}
\end{exercise}

\begin{exercise}[$\star\star\star$]\label{exo:b1:euclid:12}
मान लीजिए $u$ \omterm{def:b1:euclid:def}{यूक्लिडीय समष्टि} $E$ की ऐसी अंतःसमाकारिता है जिसका किसी
\omterm{def:b1:euclid:orthogonal}{लांबिक-प्रसामान्य} \omterm{def:b1:vspaces:free}{आधार} में आव्यूह $A$ \omterm{def:b1:euclid:orthogonal}{लांबिक} ($A^{\mathsf T}A = I$) भी है और
सममित ($A^{\mathsf T} = A$) भी।
\begin{enumerate}
  \item $A^2 = I$ दिखाइए, और (\cref{thm:b1:linmaps:projchar} के द्वारा)
  निकालिए कि $E = \ker(u - \mathrm{id}) \oplus \ker(u + \mathrm{id})$।
  \item दिखाइए कि दोनों उपसमष्टियाँ \omterm{def:b1:euclid:orthogonal}{लांबिक} हैं, अतः $u$ $F = \ker(u -
  \mathrm{id})$ के सापेक्ष \emph{\omterm{def:b1:euclid:orthogonal}{लांबिक} सममिति} है: अर्थात् $F$ में
  परावर्तन। \emph{($u(x) = x$ और $u(y) = -y$ के लिए $\langle x, y\rangle$ दो
  तरीक़ों से संगणित कीजिए।)}
  \item समतल वाली स्थिति का वर्गीकरण कीजिए: \cref{thm:b1:euclid:plane} के
  कौन-से आव्यूह सममित हैं, और तदनुरूप \omterm{def:b1:logic:map}{प्रतिचित्रण} क्या हैं?
\end{enumerate}
\end{exercise}

\section{समस्या: समतल के समदूरक और लियोनार्दो की प्रमेय}

\begin{problem}\label{pb:b1:euclid:1}
\emph{समतल का \omterm{def:b1:euclid:isometry}{समदूरक}} कोई भी ऐसा \omterm{def:b1:logic:map}{प्रतिचित्रण} $f \colon \R^2 \to \R^2$ है जो
दूरियाँ सुरक्षित रखे: सभी $x, y$ के लिए $\norm{f(x) - f(y)} = \norm{x - y}$
--- रैखिकता की कोई परिकल्पना नहीं। यह समस्या सिद्ध करती है कि ऐसे
\omterm{def:b1:logic:map}{प्रतिचित्रण} ठीक स्थानांतरण, घूर्णन, परावर्तन और \omterm{pb:b1:euclid:1}{सर्पी परावर्तन} ही हैं
(समतल-समदूरकों का वर्गीकरण), उनके संयोजन संगणित करती है, और उनके सभी
\emph{परिमित} \omterm{def:b1:structures:group}{समूह} निर्धारित करती है: यही लियोनार्दो की प्रमेय है, अर्थात्
पुष्पिका-प्रतिरूपों के पीछे का गणित। भाग II से आगे हम $\R^2$ को $\C$ के साथ
एक मान लेते हैं (\cref{ch:b1:complex}): विहित \omterm{def:b1:euclid:def}{अंतर्गुणन} $\langle z, w\rangle
= \operatorname{Re}(z\conj w)$ है और मानक \omterm{def:b1:complex:field}{मापांक}।\index{समदूरक}\index{सर्पी
परावर्तन} \index{लियोनार्दो की प्रमेय}\index{द्विफलकी समूह}

\medskip \noindent\textbf{भाग I --- हर \omterm{def:b1:euclid:isometry}{समदूरक} एफ़ीन है।}
\begin{enumerate}
  \item जाँचिए कि स्थानांतरण $t_a(x) = x + a$, \omterm{def:b1:linmaps:def}{रैखिक} \omterm{def:b1:euclid:isometry}{समदूरक}, और उनके सभी
  संयोजन \omterm{def:b1:euclid:isometry}{समदूरक} हैं, तथा यह कि \omterm{def:b1:euclid:isometry}{समदूरक} संयोजन के अंतर्गत एक \omterm{def:b1:structures:group}{समूह} बनाते हैं।
  \item मान लीजिए $f$ ऐसा \omterm{def:b1:euclid:isometry}{समदूरक} है कि $f(0) = 0$। दिखाइए कि $f$ मानक
  सुरक्षित रखता है, फिर --- ध्रुवण से, \cref{thm:b1:euclid:cs} --- कि सभी
  $x, y$ के लिए $\langle f(x), f(y)\rangle = \langle x, y\rangle$।
  \item अब भी $f(0) = 0$ के साथ: प्रश्न 2 का प्रयोग करके $\norm{f(x + y) -
  f(x) - f(y)}^2$ और $\norm{f(\lambda x) - \lambda f(x)}^2$ का प्रसार कीजिए,
  और निष्कर्ष निकालिए कि $f$ \emph{\omterm{def:b1:linmaps:def}{रैखिक}} है: $f \in O(\R^2)$।
  \item इससे निष्कर्ष निकालिए कि हर \omterm{def:b1:euclid:isometry}{समदूरक} $f$ \emph{अद्वितीय} रूप से $f =
  t_a \circ g$ लिखा जाता है, जहाँ $a = f(0)$ और $g$ कोई \omterm{def:b1:linmaps:def}{रैखिक} \omterm{def:b1:euclid:isometry}{समदूरक} है ($f$
  का \emph{\omterm{def:b1:linmaps:def}{रैखिक} भाग})।
  \item $f$ को \emph{प्रत्यक्ष} कहिए जब $\det g = 1$, और \emph{अप्रत्यक्ष}
  जब $\det g = -1$। दिखाइए कि किसी संयोजन का \omterm{def:b1:linmaps:def}{रैखिक} भाग \omterm{def:b1:linmaps:def}{रैखिक} भागों का संयोजन
  है, और उससे मिलने वाला चिह्न-नियम कहिए (प्रत्यक्ष/अप्रत्यक्ष $+1/-1$ की
  तरह जुड़ते हैं)।
\end{enumerate}

\medskip \noindent\textbf{भाग II --- चार प्रकार।} \cref{thm:b1:euclid:plane}
के द्वारा $\C$ के \omterm{def:b1:linmaps:def}{रैखिक} \omterm{def:b1:euclid:isometry}{समदूरक} $z \mapsto az$ और $z \mapsto a\conj z$ हैं,
जहाँ $\abs a = 1$; अतः समतल का हर \omterm{def:b1:euclid:isometry}{समदूरक} है
\[
f(z) = a z + b \quad (\text{प्रत्यक्ष}) \qquad\text{अथवा}\qquad
f(z) = a\conj z + b \quad (\text{अप्रत्यक्ष}), \qquad \abs a = 1 .
\]
\begin{enumerate}[resume]
  \item शब्दकोश सत्यापित कीजिए: $R_\theta$ $z \mapsto \eu^{\iu\theta}z$ है
  और $S_\theta$ $z \mapsto \eu^{\iu\theta}\conj z$ है ($1$ तथा $\iu$ दोनों
  पर जाँचिए)।
  \item (प्रत्यक्ष स्थिति) मान लीजिए $f(z) = az + b$, $\abs a = 1$। दिखाइए:
  यदि $a = 1$, तो $f$ एक स्थानांतरण है; और यदि $a \neq 1$, तो $f$ का
  अद्वितीय अचल बिंदु $z_0 = b/(1 - a)$ है और $f(z) - z_0 = a(z - z_0)$:
  अर्थात् केंद्र $z_0$ तथा कोण $\arg a$ वाला एक \emph{घूर्णन}।
  \item (अप्रत्यक्ष स्थिति) मान लीजिए $f(z) = a\conj z + b$, और $v = a\conj
  b + b$। $f \circ f = t_v$ तथा $f \circ t_v = t_v \circ f$ दिखाइए। यदि $v =
  0$: तो दिखाइए कि $z$ और $f(z)$ का मध्यबिंदु अचल बिंदु है, और यह कि $f$
  किसी रेखा में \emph{परावर्तन} है। यदि $v \neq 0$: तो दिखाइए कि $r =
  t_{-v/2}\circ f$ ऐसा परावर्तन है जिसका अक्ष $v$ के समांतर है, अतः $f =
  t_{v/2} \circ r$ एक \emph{\omterm{pb:b1:euclid:1}{सर्पी परावर्तन}} है। निष्कर्ष निकालिए: समतल का हर
  \omterm{def:b1:euclid:isometry}{समदूरक} कोई स्थानांतरण, घूर्णन, परावर्तन अथवा \omterm{pb:b1:euclid:1}{सर्पी परावर्तन} है
  (समतल-समदूरकों का वर्गीकरण)।
  \item (संयोजन) दिखाइए: $\alpha$ और $\beta$ कोणों के घूर्णनों का संयोजन
  $\alpha + \beta$ कोण का घूर्णन है ($\alpha + \beta \in 2\pi\Z$ होने पर एक
  स्थानांतरण); और दो परावर्तनों का संयोजन ऐसा घूर्णन है जिसका कोण दोनों
  अक्षों के बीच के कोण का दुगुना है (अक्षों के समांतर होने पर एक
  स्थानांतरण)।
  \item इससे निष्कर्ष निकालिए कि समतल का हर \omterm{def:b1:euclid:isometry}{समदूरक} अधिक से अधिक तीन
  परावर्तनों का संयोजन है।
\end{enumerate}

\medskip \noindent\textbf{भाग III --- तीन पहचानें।}
\begin{enumerate}[resume]
  \item $f(z) = \iu\conj z + 1 + \iu$ का पूरा वर्गीकरण कीजिए: प्रकार, अक्ष,
  सर्पण सदिश।
  \item मान लीजिए $f$ $0$ के परितः $\frac\pi2$ कोण का घूर्णन है और $g$ $1$
  के परितः $\frac\pi2$ कोण का। $g \circ f$ को $z \mapsto az + b$ के रूप में
  संगणित कीजिए और उसे पहचानिए (प्रकार, केंद्र, कोण)।
  \item मान लीजिए $r_1(z) = \conj z$ (वास्तविक अक्ष में परावर्तन) और $r_2(z)
  = \iu\conj z$ (रेखा $y = x$ में परावर्तन)। $r_2 \circ r_1$ संगणित कीजिए और
  इस उदाहरण पर प्रश्न 9 जाँचिए।
\end{enumerate}

\medskip \noindent\textbf{भाग IV --- परिमित \omterm{def:b1:structures:group}{समूह}: लियोनार्दो की प्रमेय।} मान
लीजिए $G$ समतल-समदूरकों का कोई \emph{परिमित} \omterm{def:b1:structures:group}{समूह} है।
\begin{enumerate}[resume]
  \item दिखाइए कि \omterm{def:b1:euclid:isometry}{समदूरक} भारकेंद्र सुरक्षित रखते हैं: यदि $\lambda_1 + \dots
  + \lambda_m = 1$ और $f = t_a \circ g$ ($g$ \omterm{def:b1:linmaps:def}{रैखिक}), तो $f\bigl(\sum_i
  \lambda_i x_i\bigr) = \sum_i \lambda_i f(x_i)$।
  \item किसी भी चुने हुए $x_0$ के लिए $c = \frac{1}{\abs G}\sum_{g \in G}
  g(x_0)$ रखिए। दिखाइए कि हर $h \in G$ $c$ को अचल रखता है: अर्थात् समदूरकों
  के किसी परिमित \omterm{def:b1:structures:group}{समूह} का कोई उभयनिष्ठ अचल बिंदु होता है।
  \item इससे निष्कर्ष निकालिए कि $t_{-c}$ से संयुग्मन के बाद $G \subseteq
  O(\R^2)$ मान लिया जा सकता है। मान लीजिए $G^{+} = \{g \in G : \det g =
  1\}$; दिखाइए कि या तो $G = G^{+}$, अथवा $G^{+}$ का $G$ में सूचकांक ठीक $2$
  है (कोई एकैकी आच्छादन $G^+ \to G \setminus G^+$ दिखाइए)।
  \item दिखाइए कि $c$ के परितः \emph{घूर्णनों} का कोई परिमित \omterm{def:b1:structures:group}{समूह} चक्रीय
  होता है: उसके अवयवों में से सबसे छोटे कोण $\theta_0 \in \intoo{0}{2\pi}$
  वाला घूर्णन $R_{\theta_0}$ चुनिए, और कोणों के यूक्लिडीय भाग से सिद्ध कीजिए
  कि वही जनक है; निष्कर्ष निकालिए $\theta_0 = \frac{2\pi}n$ और $G^{+} =
  \{R_{\theta_0}^{\,k}\} \cong C_n$।
  \item मान लीजिए $G \neq G^{+}$ और कोई परावर्तन $s \in G$ चुनिए। दिखाइए कि
  $G = G^{+} \cup sG^{+}$, कि हर घूर्णन $r \in G^{+}$ के लिए $s r s =
  r^{-1}$, और यह कि $sG^{+}$ के सभी $n$ अवयव परावर्तन हैं: अतः $G$ कोटि $2n$
  का \emph{\omterm{pb:b1:euclid:1}{द्विफलकी समूह}} $D_n$ है।
  \item निष्कर्ष निकालिए (\emph{लियोनार्दो की प्रमेय}): समतल के समदूरकों का
  हर परिमित \omterm{def:b1:structures:group}{समूह} चक्रीय $C_n$ है अथवा द्विफलकी $D_n$।
\end{enumerate}

\medskip \noindent\textbf{भाग V --- लाभांश, और संश्लेषण।}
\begin{enumerate}[resume]
  \item सीधे दिखाइए कि समदूरकों के किसी परिमित \omterm{def:b1:structures:group}{समूह} में तत्समक के अतिरिक्त न
  कोई स्थानांतरण हो सकता है, न कोई \omterm{pb:b1:euclid:1}{सर्पी परावर्तन} (ऐसे अवयव की घातों पर
  विचार कीजिए)।
  \item मान लीजिए $P_n$ वह सम $n$-भुज है जिसके शीर्ष $n$-वें इकाई-मूल हैं
  ($n \geq 3$)। दिखाइए कि उसका सममिति-समूह ठीक $D_n$ है: $n$ घूर्णन $z
  \mapsto \omega^k z$ और $n$ परावर्तन $z \mapsto \omega^k \conj z$, $\omega
  = \eu^{2\iu\pi/n}$, और कोई नहीं।
  \item ऐसी समतल आकृतियाँ दिखाइए जिनके सममिति-समूह क्रमशः $C_1$, $D_1$,
  $D_2$ और $C_3$ हों।
  \item $\pm1, \pm\iu$ शीर्षों वाले वर्ग के सममिति-समूह के आठों अवयव $z
  \mapsto \omega^k z$ अथवा $z \mapsto \omega^k\conj z$ रूप के प्रतिचित्रणों
  के रूप में सूचीबद्ध कीजिए, और चारों परावर्तनों में से हर एक का अक्ष दीजिए।
  \item मान लीजिए $f, g$ \emph{भिन्न} केंद्रों $c_1 \neq c_2$ के परितः एक ही
  कोण $\theta \notin 2\pi\Z$ के घूर्णन हैं। $f\circ g - g\circ f$ को
  बिंदुवार संगणित कीजिए और $f\circ g \neq g\circ f$ दिखाइए; इसके अतिरिक्त
  दिखाइए कि $(f \circ g)\circ(g\circ f)^{-1}$ एक अतुच्छ स्थानांतरण है, अतः
  $f$ और $g$ को समाहित करने वाला कोई भी \omterm{def:b1:structures:group}{समूह} अनंत है --- यह लियोनार्दो की
  प्रमेय में एकमात्र केंद्र होने की दूसरी व्याख्या है।
  \item संश्लेषण, चार वाक्यों में: कौन-से दो संरचनात्मक परिणाम स्वेच्छ
  समदूरकों को रैखिक बीजगणित तक (प्रश्न 3--4) और स्वेच्छ परिमित समूहों को
  $O(2)$ के उपसमूहों तक (प्रश्न 15) घटा देते हैं; समतल-समदूरकों की पूरी सूची
  क्या है और कौन-से अचर (प्रत्यक्ष/अप्रत्यक्ष, अचल बिंदु) चारों प्रकारों को
  अलग करते हैं; प्रश्न 9 के संयोजन-नियम परावर्तनों को हर चीज़ का जनक क्यों
  बना देते हैं; और लियोनार्दो की प्रमेय परिमित पैमाने पर क्या जोड़ती है। भाग
  II और IV में सिद्ध दोनों प्रमेयों के नाम बताइए।
\end{enumerate}
\end{problem}
